\documentclass{article}
\usepackage[UTF8]{ctex}
\usepackage{amsmath,amssymb,amsfonts}
\usepackage{enumitem}
\usepackage{geometry}
\geometry{a4paper, margin=1in}
% === 错题本排版优化（不参与题目标识） ===
\setlength{\parindent}{0pt}
\setlength{\parskip}{0.35em plus 0.10em minus 0.08em}
\setlength{\emergencystretch}{3em}
\setlist{itemsep=0.25em, topsep=0.40em, parsep=0pt}
\allowdisplaybreaks[2]
% === 错题本排版优化结束 ===
\begin{document}
\section*{电子科技大学2025年数学分析真题}
\section*{168．电子科技大学2025年数学分析试卷}

\subsection{题目1}
1．求 \(\lim _{n \rightarrow \infty} n\left[\left(1+\frac{1}{n}\right)^{n}-e\right]=\) \(\_\_\_\_\)．

\subsection{题目2}
2 ．已知 \(z=x \ln \frac{x}{y}\) ，求 \(\left.\mathrm{d}^{2} z\right|_{(1,1)}=\) \(\_\_\_\_\)．

\subsection{题目3}
3．计算 \(\oint_{L}(x y+y z+z x) \mathrm{d} s=\) \(\_\_\_\_\)，其中 \(L\) 为 \(x^{2}+y^{2}+z^{2}=a^{2}\) 与 \(x+y+z=0\) 的交线。\\
\(\_\_\_\_\)

\subsection{题目4}
4．幂级数 \(\sum_{n=1}^{\infty}\left(1+\frac{1}{2}+\cdots+\frac{1}{n}\right)^{2} x^{n}\) 的收敛域为 \(\_\_\_\_\)．

\subsection{题目5}
5．已知 \(f(x)=x(\pi-x), x \in[0, \pi]\) ，展开为余弦级数为 \(\_\_\_\_\)．

\subsection{题目6}
6．已知 \(b>a>0\) ，计算 \(\int_{0}^{1} \frac{x^{b}-x^{a}}{\ln x} \mathrm{~d} x=\) \(\_\_\_\_\)．

\subsection{题目7}
7．已知数列 \(\left\{a_{n}\right\}: a_{1}=0, a_{2 m}=\frac{a_{2 m-1}}{2}, a_{2 m+1}=\frac{1}{2}+a_{2 m}\) ，求 \(\left\{a_{n}\right\}\) 的上下极限．

\subsection{题目8}
8．求 \(\int_{0}^{\frac{\pi}{2}} \frac{f(x)}{\sqrt{x}} \mathrm{~d} x\) ，其中 \(f(x)=\int_{\sqrt{\frac{\pi}{2}}}^{\sqrt{x}} \frac{1}{1+\tan u^{2}} \mathrm{~d} u\) ．

\subsection{题目9}
9．计算 \(\lim _{n \rightarrow \infty} \frac{1}{n} \sum_{j=1}^{n} \sum_{k=1}^{n} \frac{j}{j^{2}+k^{2}}\) ．

\subsection{题目10}
10．计算 \(\int_{L} \frac{y \mathrm{~d} x-(x-1) \mathrm{d} y}{(x-1)^{2}+y^{2}}\) ，其中 \(L\) 为椭圆 \(x^{2}-4 x+4 y^{2}=0\) ，取逆时针方向。

\subsection{题目11}
11．求第二类曲面积分 \(\iint_{S} \frac{x}{x^{2}+y^{2}+z^{2}} \mathrm{~d} y \mathrm{~d} z\) ，其中 \(S\) 为圆柱面 \(x^{2}+y^{2}=1(-1 \leq z \leq 1)\) ，方向为外侧。

\subsection{题目12}
12．\(A\) 是由数码 0,1 组成的所有数列的集合，证明 \(A\) 不可数．

\subsection{题目13}
13．证明函数 \(f(x)=\sqrt{x} \ln x\) 在 \([1,+\infty)\) 上一致连续．

\subsection{题目14}
14．证明级数 \(\sum_{n=1}^{\infty} \frac{\sin (n+1) \cos (n-1)}{n^{p}}\) 在 \(p>1\) 时收敛，在 \(p \leq 1\) 时发散。

\subsection{题目15}
15．证明级数 \(\sum_{n=1}^{\infty}(-1)^{n} \frac{x^{2}+n}{n^{2}}\) 在闭区间 \([a, b]\) 上一致收敛，但对任意的 \(x\) 不绝对一致收敛．

\subsection{题目16}
16．\((x)\) 表示 \(x\) 的小数部分，已知 \(f(x)=\sum_{n=1}^{\infty} \frac{(n x)}{n^{2}}(x \in \mathbb{R})\) ．

\begin{itemize}
  \item[（1）] 求 \(f(x)\) 所有间断点．
  \item[（2）] 证明 \(f(x)\) 在有界闭区间 \([a, b]\) 上黎曼可积．
\end{itemize}

\subsection{题目17}
17．已知数列 \(\left\{a_{n}\right\}\) 满足 \(a_{n+1}>a_{n}>0\) ，且 \(\sum_{n=1}^{\infty} a_{n}\) 收敛．

\begin{itemize}
  \item[（1）] 证明： \(\lim _{n \rightarrow \infty} n a_{n}=0\) ．
  \item[（2）] 证明：\(\sum_{n=1}^{\infty} n\left(a_{n}-a_{n+1}\right)\) 收敛。
\end{itemize}
\end{document}
